\documentclass[../Main.tex]{subfiles}

\begin{document}
\section{Velocity Potential}
In Chapter 2 we saw that if at $t = 0$, $\vec{\omega} = \vec0$, then $\vec{\omega} = \vec0$ at all time. That is, irrotational flow remains irrotational.

If $\vec{\omega} = \nabla \times \vec{u} = \vec0$, then we know from the course IA Vector Calculus that:
\begin{equation*}
    \vec{u}(\vec{x}, t) = \nabla \phi(\vec{x}, t)
\end{equation*}
where $\phi$ is the \underline{velocity potential}.
\begin{remark}
    We use a positive grad here, unlike, for example, $\bdforce = - \nabla \chi$.
\end{remark}
By applying the incompressibility of the flow, we get $\nabla \cdot (\nabla \phi) = 0$. We also apply kinematic boundary conditions to get the problem:
\begin{equation}
    \begin{cases}
        \nabla^2 \phi = 0 & \vec{x} \in \Omega \\
        \frac{\partial \phi}{\partial \vec{n}} = \vec{U} \cdot \vec{n} & \vec{x} \in \partial \Omega
    \end{cases}
    \label{eqnPotentialLaplace}
\end{equation}
Therefore, the Euler equation reduces to the Laplace equation (with Neumann boundary conditions). We use the Euler equation after we have found the flow in order to get the pressure.
\section{Simple Examples}
For simple geometries, we can build solutions using separable solutions to the Laplace Equation in suitable coordinate systems.

The simplest example we could think of is uniform flow $\vec{u} = \vec{U}$. Then $\phi = \vec{U} \cdot \vec{x}$.

\begin{example}[Spherical coordinates]
    The most general axisymmetric solution to Laplace's equation is:
    \begin{equation*}
        \phi = \sum_{n=0}^{\infty} \left(A_n r^n + B_n r^{-1-n}\right) P_n(\cos(\theta))
    \end{equation*}
    where $P_n$ are the Legendre polynomials.

    Usually we only need a small number of modes.
    \begin{itemize}
        \item The $B$ term when $n = 0$ is $\phi = B / r$. This gives a \underline{3D point source}, with velocity $\vec{u} = -\frac{B}{r^2} \vec{e_r}$. This is a radial flow with constant volume flux through a sphere:
            \begin{equation*}
                \int_{B_R(\vec0)} \vec{u} \cdot \vec{n} dS = -\frac{B}{R^2} \times 4\pi R^2 = -4\pi B
            \end{equation*}
            we define this to be the volume flux $q$, so:
            \begin{equation*}
                \phi = -\frac{q}{4\pi r}
            \end{equation*}
        \item The $B$ term when $n = 1$ is:
            \begin{equation*}
                \phi = B \frac{\cos(\theta)}{r^2}
            \end{equation*}
            this is \underline{dipole flow}, where a source and a sink are very close to each other.
    \end{itemize}
\end{example}
\begin{example}[Cylindrical coordinates]
    Consider cylindrical coordinates with no $z$ dependence (equivalent to 2D polar coordinates).

    The most general solution to Laplace's Equation is:
    \begin{equation*}
        \rho = C + A_0 \log(r) + B_0 \theta + \sum_{n=1}^{\infty} (A_n r^n + B_n r^{-n}) (C_n cos(n\theta) + D_n \sin(n\theta))
    \end{equation*}
    Since this is only defined up to a constant, we disregard $C$.
    
    We can take just a few nodes to give:
    \item With only the $A_0$ term,
        \begin{equation*}
            \phi = \frac{q}{2\pi}\log(r) \implies \vec{u} = \frac{q}{2\pi r} \vec{e_r}
        \end{equation*}
        which is a \underline{2D point source}.
    \item With only the $B_0$ term,
        \begin{equation*}
            \phi = \frac{\kappa}{2\pi}\theta \implies \vec{u} = \frac{\kappa}{2\pi r} \vec{e_\theta}
        \end{equation*}
        then this is circular motion around the origin. The vorticity is zero except at $0$ (vorticity is a delta function). Therefore this flow is invalid if we include the origin. 
    \item $\phi = Ur\cos(\theta)$ is a solution. %TODO: How?
    \item The \underline{2D dipole} is given by:
        \begin{equation*}
            \phi = A \frac{\cos(\theta)}{r}
        \end{equation*}
\end{example}
\section{Uniform Flow Past a Sphere or Cylinder}
\subsection{Flow Past a Sphere}
\begin{figure}
    \centering
    \begin{tikzpicture}[scale=1]
    
    \end{tikzpicture}
    \caption{Diagram of flow past a sphere}
    \label{figFlowSphere}
\end{figure}
Consider figure~\ref{figSphere}. We want to solve $\nabla^2 \phi = 0$ with the boundary conditions:
\begin{align*}
    \vec{u} &\to U \vec{e_z} \text{ far from the sphere} \\
    \vec{u} \cdot \vec{n} &= \vec{U} \cdot \vec{n} \text{ (KBC)}
\end{align*}
in terms of the velocity potential $\phi$,
\begin{align*}
    \phi(r, \theta) &\to U r \cos(\theta) \text{ as $r\to\infty$} \\
    \frac{\partial \phi}{\partial r}(a, \theta) &= 0 
\end{align*}
Then we try a solution of the form $\phi = f(r) P_1(\cos(\theta)) = f(r) \cos(\theta)$. We recall the general solution to the Laplace equation and suppose that $f(r)$ is of the form:
\begin{equation*}
    f(r) = Ar + B r^{-2}
\end{equation*}
We know that this solves the Laplace equation, so it only remains to check the boundary conditions. Flow at infinity gives $U = A$, and the kinematic boundary condition gives:
\begin{align*}
    \frac{\partial \phi}{\partial r} &= \left[U - \frac{2B}{r^3}\right]\cos(\theta) \\
    \left.\frac{\partial \phi}{\partial r}\right|_{r = a} &= \left[U - \frac{2B}{a^3}\right]\cos(\theta) = 0 \\
    \therefore 2B &= Ua^3 \\
    \therefore B &= \frac{Ua^3}{2}
\end{align*}
And so our solution for $\phi$ is:
\begin{equation*}
    \phi = U \cos(\theta) \left[r + \frac{a^3}{2r^2}\right]
\end{equation*}
which gives a solution for $\vec{u}$ by taking a gradient:
\begin{equation*}
    \vec{u} = U \cos(\theta) \left(1 - \frac{a^3}{r^3}\right)\vec{e_r} - U \sin(\theta) \left(1 + \frac{a^3}{2r^3}\right)\vec{e_\theta}
\end{equation*}
\subsection{Uniform Flow Past a Cylinder with Circulation}
We consider cylindrical polar coordinates and solve $\nabla^2 \phi = 0$. The boundary conditions are:
\begin{align*}
    \phi &\xrightarrow[r \to \infty]{} Ur\cos(\theta) \\
    \frac{\partial \phi}{\partial r}(a, \theta) = 0 \\
    \oint_{B_r(\vec0)} \vec{u} \cdot \vec{dx} &= \kappa
\end{align*}
We cannot remove the circulation flow in this case, because we let the cylinder have infinite length along the $z$ axis.

The solution is:
\begin{align*}
    \phi &= \frac{\kappa}{2\pi}\theta + U\cos(\theta) \left[r + \frac{a^2}{r}\right] \\
    \therefore \vec{u} &= U\cos\left(\theta\right) (1 - \frac{a^2}{r^2}) \vec{e_r} - U\sin(\theta) \left(1 + \frac{a^2}{r^2}\right)\vec{e_\theta} + \frac{\kappa}{2\pi r} \vec{e_\theta}
\end{align*}
Importantly, we only have symmetry along the $x$ axis when $\kappa = 0$. See figure~\ref{figFlowCylinderComparison} for an illustration of this. An important consequence is that there will be a difference in pressure on either side of the cylinder, and so the fluid exerts a force on the cylinder.
\section{Pressure in a Potential Flow}
The Euler equation for pressure is:
\begin{equation*}
    \rho \frac{D \vec{u}}{D t} = - \nabla p + \bdforce
\end{equation*}
which is possibly nonlinear in $p$, even when we have found $\vec{u}$.

Now assume: $\bdforce = - \nabla \chi$, so the body forces are conservative, and $\rho$ is constant.

\begin{align*}
    0 &= \nabla p + \nabla \chi + \rho \left(\frac{\partial \vec{u}}{\partial t} + (\vec{u} \cdot \nabla)\vec{u}\right) \\
    0 &= \nabla p + \nabla \chi + \rho \frac{\partial (\nabla \phi)}{\partial t} + \rho \nabla \left(\frac12 |\vec{u}|^2\right) - \vec{u} \times \underbrace{\vec{\omega}}_{0} \\
    0 &= \nabla p + \nabla \chi + \nabla\left(\rho\frac{\partial \phi}{\partial t}\right) + \nabla \left(\rho\frac12 |\vec{u}|^2\right) \\
    0 &= \nabla\left( p + \chi + \rho\frac{\partial \phi}{\partial t} + \frac12\rho |\vec{u}|^2\right)
\end{align*}
and this is equivalent to:
\begin{equation}
    \rho \frac{\partial \phi}{\partial t} + \frac12 \rho |\nabla \phi|^2 + p + \chi = f(t)
    \label{eqnBernoulliUnsteady}
\end{equation}
where $f$ is any function of time. Importantly, it does not depend on space.
This has key differences from \eqnref{eqnBernoulliSteady}:
\begin{enumerate}
    \item \eqnref{eqnBernoulliSteady} requires steady flow, but permits vorticity, \eqnref{eqnBernoulliUnsteady} requires irrotational flow, but permits time dependence;
    \item both require conservative forces
    \item \eqnref{eqnBernoulliSteady} is just a consequence of the Euler momentum equation~\ref{eqnEulerMomDiv}, because we took a scalar product; \eqnref{eqnBernoulliUnsteady} is equivalent to the Euler equation under the given assumptions.
\end{enumerate}
\end{document}